\chapter{Materiais}\label{ch:materials}

Este capítulo descreve a área de estudo, os sete segmentos de uso da terra, a plataforma de computação
e o inventário de dados. As tabelas completas de fontes de dados estão no Anexo~\ref{anx:data}.

\section{Área de estudo}

A área de interesse compreende o estado de Goiás e o Distrito Federal (GO~e~DF), na região Centro-Oeste
do Brasil, inteiramente inseridos no bioma Cerrado. O Cerrado é uma savana tropical de clima fortemente
sazonal (verão úmido e inverno seco). Seus solos, predominantemente Latossolos, são profundos,
intemperizados, naturalmente ácidos e de baixa fertilidade, mas podem se tornar altamente produtivos
após a calagem e a correção de fósforo, potássio, enxofre e micronutrientes
\citep[p.~617]{yamada_cerrado_2005}. Nas últimas décadas, o bioma se tornou uma das principais regiões
produtoras de grãos, fibras e carne do Brasil. É precisamente isso que torna relevante a questão do
potencial \emph{versus} uso atual (QP2).

A geometria da área de interesse é derivada da Malha Municipal Digital oficial do IBGE (edição 2025,
SIRGAS~2000/EPSG:4674, 246 municípios de Goiás mais o Distrito Federal), dissolvida em um polígono único
que se estende aproximadamente da longitude $-53{,}25^\circ$ a $-45{,}90^\circ$ e da latitude
$-19{,}50^\circ$ a $-12{,}39^\circ$. Na fase de uso atual da terra, as 247 unidades municipais dessa
malha servem de base para a agregação e a apresentação dos resultados por município.

\section{Os sete segmentos de uso da terra}

O atlas avalia sete segmentos, escolhidos de forma a cobrir produção agrícola, fibras, proteína animal,
energia e conservação, para que a análise possa expressar compensações reais (\emph{trade-offs}) entre
segmentos. A Tabela~\ref{tab:segments} caracteriza cada segmento em termos agronômicos; a literatura de
mapeamento multicritério por segmento é revisada na Seção~\ref{sec:sota-casos}.

\begin{table}[htbp]
\centering
\caption{Os sete segmentos de uso da terra.}\label{tab:segments}
\begin{tabular}{@{}cl>{\raggedright\arraybackslash}p{9.5cm}@{}}
\toprule
\# & Segmento & Definição e escopo \\
\midrule
1 & \textbf{Soja} & Soja mecanizada de sequeiro, principal cultura comercial da região. Seu desempenho
  responde fortemente a covariáveis edafoclimáticas, como disponibilidade hídrica do solo, temperatura
  e fotoperíodo, como mostram ensaios multiambiente em Mato Grosso
  \citep[pp.~1--3]{corbellini_geographical_2024}. Projeções para a área de influência do Arco Norte
  indicam recuo de 12,6\% a 34,8\% da área de soja em Goiás em 2041--2060
  \citep[p.~11]{souza_climate-driven_2026}. \\
2 & \textbf{Cana-de-açúcar} & Cana-de-açúcar mecanizada para açúcar e etanol, com maturação na estação
  seca. As faixas de temperatura e precipitação partem do zoneamento de aptidão climática da cultura no
  Brasil \citep[p.~782]{aparecido_climate_2021} e foram recalibradas às distribuições observadas em
  GO~e~DF (Anexo~\ref{anx:spec}). \\
3 & \textbf{Outras culturas anuais} & Milho, algodão e sorgo, incluindo a safrinha (segunda safra).
  Os critérios e as faixas de classe seguem avaliações SIG-AHP publicadas para cereais de sequeiro
  \citep[pp.~6--8]{tadesse_land_2020}. \\
4 & \textbf{Piscicultura} & Aquicultura em viveiros escavados e reservatórios, com espécies de água quente,
  como a tilápia. Os critérios (temperatura, textura do solo, declividade e acesso à água) seguem estudos
  de aquicultura em viveiros em outras regiões
  \parencites{ssegane_geospatial_2012}[Tabela~1]{assefa_gis_2018}; a temperatura do ar é usada como
  aproximação da temperatura da água. \\
5 & \textbf{Pecuária} & Gado de corte em pastagem cultivada, que cobre cerca de 28\% do bioma Cerrado e
  abrange um amplo espectro de estados de intensificação e degradação
  \parencites[p.~1219]{vieira_classification_2022}{molossi_improve_2020}. \\
6 & \textbf{Conservação} & Conservação de floresta e Cerrado: valor dos serviços da vegetação nativa,
  da água e da biodiversidade, em consonância com as compensações entre prioridade de restauração e
  aptidão agrícola mapeadas para o bioma \citep{schuler_spatial_2022}. \\
7 & \textbf{Geração fotovoltaica} & Geração fotovoltaica centralizada, de grande porte. A escolha do local
  considera irradiância, índice de claridade, orientação da vertente, temperatura e declividade, sobre a
  máscara de cobertura da terra, em linha com os fatores usuais em análises multicritério para usinas
  fotovoltaicas \citep[Tabela~2]{rane_gis-based_2024}. \\
\bottomrule
\end{tabular}
\end{table}

\section{Plataforma de computação e filosofia de custo}\label{sec:cost}

A análise é executada inteiramente no servidor do Google Earth Engine, plataforma em nuvem que
integra um catálogo de dados de observação da Terra a um serviço de computação paralela
\citep[pp.~18--19]{gorelick_google_2017}. Utiliza-se aqui a modalidade gratuita, para uso não comercial e
de pesquisa. O objetivo dessa abordagem é a reprodutibilidade com baixo custo marginal: sem
\emph{hardware} dedicado e sem armazenamento local de imagens de satélite. A máquina local é usada apenas
para geração de gráficos, junção de tabelas e operações simples. Os produtos intermediários são
armazenados na plataforma como ativos (\emph{assets}) matriciais persistentes, para que as etapas
seguintes não recalculem os grafos anteriores. Para que o processamento permaneça dentro da cota de
uso, as amostras de agrupamento e classificação são limitadas à ordem de $10^5$ pontos.

\begin{table}[htbp]
\centering
\caption{Grade de análise e períodos de referência.}\label{tab:grid-periods}
\begin{tabular}{@{}>{\raggedright\arraybackslash}p{3.2cm}>{\raggedright\arraybackslash}p{3.8cm}>{\raggedright\arraybackslash}p{6.3cm}@{}}
\toprule
Convenção & Valor & Justificativa \\
\midrule
Resolução de análise & 250~m & Corresponde à grade nativa das camadas de solo (OpenLandMap); situa-se entre a resolução fina do terreno e a resolução mais baixa do clima. \\
Sistema de referência de coordenadas & EPSG:4326 & A área das células vem da função planimétrica de área de pixel, nunca de um campo armazenado. \\
Normal climatológica & 1991--2020 & Referência única para os deltas de mudança climática. \\
Registro da vegetação & 2001--2020 & Indicador de produtividade (MOD17). \\
Janelas futuras & 2031--2050; 2051--2070 & Duas janelas de projeção. \\
Cenários futuros & SSP2-4.5; SSP5-8.5 & Cenários de emissão moderada e alta. \\
\bottomrule
\end{tabular}
\end{table}

\section{Fontes de dados}

Com uma única exceção, todos os conjuntos de dados vêm do catálogo público da própria
plataforma; o Anexo~\ref{anx:data} registra, fonte a fonte, a resolução nativa e os fatores de
conversão para unidades físicas. A exceção é a malha municipal do IBGE, distribuída fora do catálogo:
ela é lida de um arquivo local e convertida em geometria na memória a cada execução, sem nunca ser
exportada como ativo persistente na plataforma.

As variáveis \textbf{biofísicas e agroclimáticas} (Tabela~\ref{tab:src-biophys} do Anexo) abrangem:
o relevo (SRTM); a hidrografia (HydroSHEDS); o clima e o regime hídrico, pelo conjunto TerraClimate
(mensal, $\tfrac{1}{24}^\circ$, $\approx$4,6~km) \citep{abatzoglou_terraclimate_2018}; os solos, pelas
camadas OpenLandMap SOL v0.2 a 250~m (média das profundidades de 0, 10 e 30~cm); e a água superficial
(JRC Global Surface Water) \citep{pekel_high-resolution_2016}. Completam o conjunto a cobertura da terra
do MapBiomas (Coleção~10), usada apenas como máscara; as áreas protegidas (WDPA); o carbono da biomassa
aérea \citep[p.~1]{spawn_harmonized_2020}; e o clima futuro do conjunto NASA NEX-GDDP-CMIP6
\citep{thrasher_nasa_2022}.

As entradas de \textbf{uso atual e validação} (Tabela~\ref{tab:src-realized} do Anexo) são as 247
unidades municipais da Malha Municipal Digital do IBGE (edição 2025), as classificações anuais de uso da
terra do MapBiomas \citep{souza_reconstructing_2020} e os produtos de produtividade primária líquida e
bruta (NPP e GPP, do inglês \emph{net} e \emph{gross primary productivity}) do MODIS MOD17
\citep{zhao_improvements_2005}. O uso atual da terra e o indicador de produtividade servem estritamente
para mascaramento, validação e caracterização descritiva; nunca entram no cálculo da aptidão nem do
agrupamento.

\section{Cenários futuros do clima (CMIP6)}

A projeção usa o conjunto NEX-GDDP-CMIP6 da NASA, de resolução temporal diária, em particular os campos de
precipitação e de temperatura máxima e mínima do ar próximo à superfície. Usa-se um conjunto de cinco
modelos de circulação geral (GCM, do inglês \emph{general circulation models}): ACCESS-CM2, EC-Earth3,
GFDL-ESM4, MPI-ESM1-2-HR e MRI-ESM2-0. Esse subconjunto é de uso corrente, e todos os cinco modelos
têm dados para os dois cenários de interesse.

\section{Ambiente computacional}

A implementação é orientada a configuração. A lógica analítica é modular, e todo parâmetro (limiares de
pertinência \emph{fuzzy}, identificadores de conjuntos de dados, fatores de conversão de unidades, âncoras
ordinais da matriz AHP e o expoente da adaptação regional) é declarado como valor versionado, e não
embutido no código. Assim, a análise pode ser auditada e reexecutada apenas a partir da especificação
declarada. Os \emph{pesos} dos fatores são a exceção deliberada: eles não são declarados, e sim
\textbf{derivados} das âncoras pela cadeia descrita na Seção~\ref{sec:ahp}. Além da integração com o
Google Earth Engine e da orquestração das etapas de análise, o código Python usa apenas bibliotecas
científicas comuns nas tarefas locais. Nenhuma técnica de aprendizado profundo é utilizada.

O código-fonte completo (configurações, módulos de análise e ferramentas de exportação) está disponível
publicamente no repositório \texttt{goias\_suitability} \citep{da_silva_goias_suitability_2026}. O atlas
interativo construído a partir desses resultados está publicado em
\url{https://ahxsahxs.github.io/goias_suitability}.
